\chapter{Chapter01 — Parity Basis}
%%% generated by the blueprint pipeline from FABL.Chapter01.ParityBasis
\section{Overview}

Declarations of the Lean module \texttt{FABL.Chapter01.ParityBasis}.

\section{Declarations}

\begin{definition}[Uniform inner product]
\lean{FABL.uniformInner}
\statementsource{lmn-constant-depth-fourier}{pdf-63759bd909d5-p003-b006}
\statementsource{Analysis-of-Boolean-Functions-by-Ryan-ODonnell}{pdf-ffbe9c7b3306-p024-b001}
\statementsource{ghsz-combinatorial-list-decoding}{pdf-6bd33ae73ade-p005-b002}
\label{FABL.uniformInner}
\leanok
For a finite set $\Omega$ and real-valued functions $f,g$ on $\Omega$, their uniform
inner product is $\langle f,g\rangle_u=\frac{1}{|\Omega|}\sum_{x\in\Omega}f(x)g(x)$.
\end{definition}

\begin{definition}[Uniform $L^p$ expression]
\lean{FABL.uniformLpNorm}
\statementsource{lmn-constant-depth-fourier}{pdf-63759bd909d5-p004-b001}
\statementsource{Analysis-of-Boolean-Functions-by-Ryan-ODonnell}{pdf-ffbe9c7b3306-p024-b001}
\label{FABL.uniformLpNorm}
\leanok
For a finite set $\Omega$, a real number $p$, and a real-valued function $f$ on
$\Omega$, define
$\|f\|_{p,u}=\left(\frac{1}{|\Omega|}\sum_{x\in\Omega}|f(x)|^p\right)^{1/p}$, with real
powers understood as real powers.
\end{definition}

\begin{definition}[Uniform probability mass function]
\lean{FABL.uniformPMF}
\statementsource{Vad12}{pdf-bf2ce5422088-p011-b003}
\statementsource{Analysis-of-Boolean-Functions-by-Ryan-ODonnell}{pdf-ffbe9c7b3306-p024-b002}
\label{FABL.uniformPMF}
\leanok
For a finite nonempty set $\Omega$, the uniform probability mass function assigns
probability $1/|\Omega|$ to each element of $\Omega$.
\end{definition}

\begin{theorem}[Uniform integral equals expectation]
\lean{FABL.integral_uniformPMF_eq_expect}
\label{FABL.integral_uniformPMF_eq_expect}
\leanok
\uses{FABL.uniformPMF}
\difficulty{3}
Let $\Omega$ be a finite nonempty set equipped with a measurable space in which every
singleton is measurable, and let $f:\Omega\to\mathbb{R}$. The integral of $f$ with
respect to the measure induced by the uniform probability mass function on $\Omega$
equals the expectation of $f$ under that uniform distribution.
\end{theorem}

\begin{theorem}[Distinctness of Walsh characters]
\lean{FABL.binaryWalshCharacter_injective}
\label{FABL.binaryWalshCharacter_injective}
\leanok
\uses{FABL.binarySign, FABL.coordinateSum, FABL.χ}
\difficulty{5}
For every $n\ge 0$, distinct subsets $S,T\subseteq\{0,\ldots,n-1\}$ define distinct
Walsh characters on $\mathbb F_2^n$; equivalently, the map sending each subset to its
character $x\mapsto(-1)^{\sum_{i\in S}x_i}$ is injective.
\end{theorem}

\begin{definition}[Binary Walsh basis]
\lean{FABL.binaryWalshBasis}
\statementsource{lmn-constant-depth-fourier}{
  pdf-63759bd909d5-p004-b002,
  pdf-63759bd909d5-p004-b003
}
\statementsource{Analysis-of-Boolean-Functions-by-Ryan-ODonnell}{
  pdf-ffbe9c7b3306-p023-b001,
  pdf-ffbe9c7b3306-p023-b006,
  pdf-ffbe9c7b3306-p209-b005
}
\statementsource{mossel-neeman-gaussian-noise-stability}{pdf-b53c1a5aa27d-p034-b004}
\statementsource{ghsz-combinatorial-list-decoding}{pdf-6bd33ae73ade-p005-b002}
\label{FABL.binaryWalshBasis}
\leanok
\uses{FABL.binaryWalshCharacter_injective, FABL.χ, FABL.𝔽₂}
For each natural number $n$, the Walsh characters indexed by subsets
$S\subseteq\{0,\ldots,n-1\}$ form a basis over $\mathbb R$ for the space of real-valued
functions on the binary cube $\{0,1\}^n$, where the character indexed by $S$ is
$x\mapsto(-1)^{\sum_{i\in S}x_i}$.
\end{definition}

\begin{definition}[Walsh basis on the sign cube]
\lean{FABL.walshBasis}
\statementsource{lmn-constant-depth-fourier}{
  pdf-63759bd909d5-p004-b002,
  pdf-63759bd909d5-p004-b003
}
\statementsource{Analysis-of-Boolean-Functions-by-Ryan-ODonnell}{
  pdf-ffbe9c7b3306-p023-b004,
  pdf-ffbe9c7b3306-p023-b006,
  pdf-ffbe9c7b3306-p209-b005,
  pdf-ffbe9c7b3306-p210-b004
}
\statementsource{mossel-neeman-gaussian-noise-stability}{pdf-b53c1a5aa27d-p034-b004}
\statementsource{ghsz-combinatorial-list-decoding}{pdf-6bd33ae73ade-p005-b002}
\label{FABL.walshBasis}
\leanok
\uses{FABL.monomial, FABL.signMonomialChar, FABL.signMonomialChar_injective}
For each natural number $n$, the Walsh characters indexed by subsets
$S\subseteq\{0,\ldots,n-1\}$ form a basis over $\mathbb R$ for the space of real-valued
functions on the sign cube $\{-1,1\}^n$, where the character indexed by $S$ is
$x\mapsto\prod_{i\in S}x_i$.
\end{definition}

\begin{theorem}[Parity orthonormal basis]
\lean{FABL.parity_orthonormal_basis}
\statementsource{Analysis-of-Boolean-Functions-by-Ryan-ODonnell}{
  pdf-ffbe9c7b3306-p024-b003,
  pdf-ffbe9c7b3306-p024-b004,
  pdf-ffbe9c7b3306-p024-b005,
  pdf-ffbe9c7b3306-p024-b006,
  pdf-ffbe9c7b3306-p024-b007,
  pdf-ffbe9c7b3306-p025-b001
}
\proofsource{ghsz-combinatorial-list-decoding}{pdf-6bd33ae73ade-p005-b002}
\label{FABL.parity_orthonormal_basis}
\leanok
\uses{FABL.monomial, FABL.signMonomialChar, FABL.signMonomialChar_injective, FABL.uniformInner, FABL.walshBasis}
\proofstep
  {FABL.Sign}
  {context}
  {ghsz-combinatorial-list-decoding}
  {pdf-6bd33ae73ade-p005-b002}
\proofstep
  {FABL.signValue}
  {context}
  {ghsz-combinatorial-list-decoding}
  {pdf-6bd33ae73ade-p005-b002}
\proofstep
  {FABL.signValue_one}
  {context}
  {ghsz-combinatorial-list-decoding}
  {pdf-6bd33ae73ade-p005-b002}
\proofstep
  {FABL.signValue_neg_one}
  {context}
  {ghsz-combinatorial-list-decoding}
  {pdf-6bd33ae73ade-p005-b002}
\proofstep
  {FABL.monomial}
  {context}
  {ghsz-combinatorial-list-decoding}
  {pdf-6bd33ae73ade-p005-b002}
\proofstep
  {FABL.signMonomialChar}
  {context}
  {ghsz-combinatorial-list-decoding}
  {pdf-6bd33ae73ade-p005-b002}
\proofstep
  {FABL.signMonomialChar_injective}
  {context}
  {ghsz-combinatorial-list-decoding}
  {pdf-6bd33ae73ade-p005-b002}
\proofstep
  {FABL.fourierCoeff}
  {context}
  {ghsz-combinatorial-list-decoding}
  {pdf-6bd33ae73ade-p005-b002}
\proofstep
  {FABL.fourierSupport}
  {context}
  {ghsz-combinatorial-list-decoding}
  {pdf-6bd33ae73ade-p005-b002}
\proofstep
  {FABL.mem_fourierSupport}
  {context}
  {ghsz-combinatorial-list-decoding}
  {pdf-6bd33ae73ade-p005-b002}
\proofstep
  {FABL.granularityFrequencyWithin}
  {context}
  {ghsz-combinatorial-list-decoding}
  {pdf-6bd33ae73ade-p005-b002}
\proofstep
  {FABL.mem_granularityFrequencyWithin}
  {context}
  {ghsz-combinatorial-list-decoding}
  {pdf-6bd33ae73ade-p005-b002}
\proofstep
  {FABL.uniformInner}
  {context}
  {ghsz-combinatorial-list-decoding}
  {pdf-6bd33ae73ade-p005-b002}
\proofstep
  {FABL.walshBasis}
  {context}
  {ghsz-combinatorial-list-decoding}
  {pdf-6bd33ae73ade-p005-b002}
\difficulty{4}

% dataset-pipeline: stub -- prose pending build_informal_statements.py
Statement pending regeneration.
\end{theorem}

\begin{theorem}[Product of monomials and symmetric difference]
\lean{FABL.monomial_mul_monomial}
\proofsource{Analysis-of-Boolean-Functions-by-Ryan-ODonnell}{
  pdf-ffbe9c7b3306-p024-b005,
  pdf-ffbe9c7b3306-p024-b006,
  pdf-ffbe9c7b3306-p229-b002
}
\label{FABL.monomial_mul_monomial}
\leanok
\uses{FABL.monomial, FABL.signValue, FABL.signValue_eq_neg_one_or_one}
\proofstep
  {FABL.Sign}
  {context}
  {Analysis-of-Boolean-Functions-by-Ryan-ODonnell}
  {pdf-ffbe9c7b3306-p014-b002,
    pdf-ffbe9c7b3306-p022-b004}
\proofstep
  {FABL.signValue}
  {context}
  {Analysis-of-Boolean-Functions-by-Ryan-ODonnell}
  {pdf-ffbe9c7b3306-p014-b002,
    pdf-ffbe9c7b3306-p022-b004}
\proofstep
  {FABL.signValue_one}
  {context}
  {Analysis-of-Boolean-Functions-by-Ryan-ODonnell}
  {pdf-ffbe9c7b3306-p014-b002}
\proofstep
  {FABL.signValue_neg_one}
  {context}
  {Analysis-of-Boolean-Functions-by-Ryan-ODonnell}
  {pdf-ffbe9c7b3306-p014-b002}
\proofstep
  {FABL.signValue_eq_neg_one_or_one}
  {context}
  {Analysis-of-Boolean-Functions-by-Ryan-ODonnell}
  {pdf-ffbe9c7b3306-p014-b002}
\proofstep
  {FABL.monomial}
  {direct}
  {Analysis-of-Boolean-Functions-by-Ryan-ODonnell}
  {pdf-ffbe9c7b3306-p022-b001,
    pdf-ffbe9c7b3306-p022-b004}
\proofstep
  {FABL.fourierCoeff}
  {direct}
  {Analysis-of-Boolean-Functions-by-Ryan-ODonnell}
  {pdf-ffbe9c7b3306-p025-b004}
\proofstep
  {FABL.fourierSupport}
  {context}
  {Analysis-of-Boolean-Functions-by-Ryan-ODonnell}
  {pdf-ffbe9c7b3306-p022-b002,
    pdf-ffbe9c7b3306-p018-b003}
\proofstep
  {FABL.mem_fourierSupport}
  {context}
  {Analysis-of-Boolean-Functions-by-Ryan-ODonnell}
  {pdf-ffbe9c7b3306-p022-b002}
\proofstep
  {FABL.granularityFrequencyWithin}
  {context}
  {Analysis-of-Boolean-Functions-by-Ryan-ODonnell}
  {pdf-ffbe9c7b3306-p075-b005,
    pdf-ffbe9c7b3306-p076-b001}
\proofstep
  {FABL.mem_granularityFrequencyWithin}
  {context}
  {Analysis-of-Boolean-Functions-by-Ryan-ODonnell}
  {pdf-ffbe9c7b3306-p075-b005,
    pdf-ffbe9c7b3306-p076-b001}
\difficulty{5}
Let $n$ be a natural number, let $S,T$ be subsets of the $n$ coordinate indices, and let
$x$ be a vector of $n$ signs. The product of the monomials indexed by $S$ and $T$,
evaluated at $x$, equals the monomial indexed by their symmetric difference, evaluated
at $x$.
\end{theorem}

\begin{theorem}[Mean of a sign monomial]
\lean{FABL.expect_monomial}
\proofsource{ghsz-combinatorial-list-decoding}{pdf-6bd33ae73ade-p005-b002}
\label{FABL.expect_monomial}
\leanok
\uses{FABL.monomial, FABL.signMonomialChar, FABL.signMonomialChar_injective}
\proofstep
  {FABL.Sign}
  {context}
  {ghsz-combinatorial-list-decoding}
  {pdf-6bd33ae73ade-p005-b002}
\proofstep
  {FABL.signValue}
  {context}
  {ghsz-combinatorial-list-decoding}
  {pdf-6bd33ae73ade-p005-b002}
\proofstep
  {FABL.signValue_one}
  {context}
  {ghsz-combinatorial-list-decoding}
  {pdf-6bd33ae73ade-p005-b002}
\proofstep
  {FABL.signValue_neg_one}
  {context}
  {ghsz-combinatorial-list-decoding}
  {pdf-6bd33ae73ade-p005-b002}
\proofstep
  {FABL.monomial}
  {context}
  {ghsz-combinatorial-list-decoding}
  {pdf-6bd33ae73ade-p005-b002}
\proofstep
  {FABL.signMonomialChar}
  {context}
  {ghsz-combinatorial-list-decoding}
  {pdf-6bd33ae73ade-p005-b002}
\proofstep
  {FABL.signMonomialChar_injective}
  {context}
  {ghsz-combinatorial-list-decoding}
  {pdf-6bd33ae73ade-p005-b002}
\proofstep
  {FABL.fourierCoeff}
  {direct}
  {ghsz-combinatorial-list-decoding}
  {pdf-6bd33ae73ade-p005-b002}
\proofstep
  {FABL.fourierSupport}
  {context}
  {ghsz-combinatorial-list-decoding}
  {pdf-6bd33ae73ade-p005-b002}
\proofstep
  {FABL.mem_fourierSupport}
  {context}
  {ghsz-combinatorial-list-decoding}
  {pdf-6bd33ae73ade-p005-b002}
\proofstep
  {FABL.granularityFrequencyWithin}
  {context}
  {ghsz-combinatorial-list-decoding}
  {pdf-6bd33ae73ade-p005-b002}
\proofstep
  {FABL.mem_granularityFrequencyWithin}
  {context}
  {ghsz-combinatorial-list-decoding}
  {pdf-6bd33ae73ade-p005-b002}
\difficulty{4}
Let $n$ be a nonnegative integer and let $S$ be a subset of the $n$ coordinates. The
uniform average of the sign monomial indexed by $S$ is $1$ if $S$ is empty and $0$
otherwise.
\end{theorem}
